% Part: first-order-logic
% Chapter: beyond
% Section: many-sorted-logic

\documentclass[../../../include/open-logic-section]{subfiles}

\begin{document}

\olfileid[hi]{fol}{byd}{msl}

\olsection{बहु-वर्गीय तर्कशास्त्र}

प्रथम-कोटि तर्कशास्त्र में चर और परिमाणक एक ही !!{domain} पर विचरते हैं।
परंतु प्रायः अनेक (परस्पर असंयुक्त) !!{domain}s को रखना उपयोगी होता है: जैसे,
संख्याओं का एक !!a{domain}, ज्यामितीय वस्तुओं का एक !!a{domain}, संख्याओं से
संख्याओं में जाने वाले फलनों का एक !!a{domain}, आबेली समूहों का एक
!!a{domain}, इत्यादि।

बहु-वर्गीय तर्कशास्त्र ऐसा ढाँचा देता है। हम ``वर्गों'' की एक सूची से आरंभ
करते हैं---किसी वस्तु का ``वर्ग'' बताता है कि उसे किस ``!!{domain}'' में
होना है। फिर प्रत्येक वर्ग के लिए !!{variable}s और परिमाणकों का एक-एक भंडार,
तथा (सामान्यतः) प्रत्येक वर्ग के लिए एक !!{identity} होता है। फलनों और संबंधों को भी उन
वस्तुओं के वर्गों के अनुसार ``प्रकारित'' किया जाता है जिन्हें वे कोणांक के
रूप में ले सकते हैं। शेष नियम प्रथम-कोटि तर्कशास्त्र के सामान्य नियम ही
रहते हैं, और परिमाणक-नियमों का एक-एक रूप प्रत्येक वर्ग के लिए दोहराया जाता
है।

उदाहरणतः अंतरराष्ट्रीय संबंधों के अध्ययन के लिए हम वस्तुओं के दो वर्गों---
फ्रांसीसी नागरिकों और जर्मन नागरिकों---वाली भाषा चुन सकते हैं। पहले वर्ग की
वस्तुओं के लिए एक एकपदी संबंध ``मदिरा पीता है'', दूसरे वर्ग की वस्तुओं के
लिए एक अन्य एकपदी संबंध ``वुर्स्ट खाता है'', और दो कोणांक लेने वाला द्विपदी
संबंध ``बहुराष्ट्रीय विवाहित युगल बनाते हैं'' हो सकता है; इसमें पहला कोणांक
पहले वर्ग का और दूसरा कोणांक दूसरे वर्ग का होगा। यदि $a$, $b$, $c$ चर
फ्रांसीसी नागरिकों पर और $x$, $y$, $z$ चर जर्मन नागरिकों पर विचरते हैं, तो
\[
\lforall[a][\lforall x][(\Atom{\Obj{MarriedTo}}{a,x} \lif
(\Atom{\Obj{DrinksWine}}{a} \lor \lnot \Atom{\Obj{EatsWurst}}{x})]]
\]
कहता है कि यदि कोई फ्रांसीसी व्यक्ति किसी जर्मन से विवाहित है, तो या तो वह
फ्रांसीसी मदिरा पीता है अथवा वह जर्मन वुर्स्ट नहीं खाता।

बहु-वर्गीय तर्कशास्त्र को स्वाभाविक रीति से प्रथम-कोटि तर्कशास्त्र में
अंतःस्थापित किया जा सकता है। इसके लिए बहु-वर्गीय !!{domain}s की सभी वस्तुओं
को एक प्रथम-कोटि !!{domain} में मिला देते हैं, वर्गों का हिसाब रखने के लिए
एकपदी !!{predicate}s का प्रयोग करते हैं और परिमाणकों को सापेक्ष बनाते हैं।
उदाहरणतः ऊपर के उदाहरण के अनुरूप प्रथम-कोटि भाषा में बताए गए अन्य संबंधों
के साथ एकपदी !!{predicate}s ``$\Obj{German}$'' और ``$\Obj{French}$'' होंगे,
और वर्ग-संबंधी प्रतिबंध हटा दिए जाएँगे। वर्गित परिमाणक $\lforall[x][!A]$,
जहाँ $x$ जर्मन वर्ग का !!a{variable} है, का अनुवाद होगा
\[
\lforall[x][(\Atom{\Obj{German}}{x} \lif !A)].
\]
हमें ऐसे अभिगृहीत जोड़ने होंगे जो सुनिश्चित करें कि वर्ग अलग-अलग हैं---जैसे,
$\lforall[x][\lnot (\Atom{\Obj{German}}{x} \land
  \Atom{\Obj{French}}{x})]$---और ऐसे अभिगृहीत भी जो प्रत्याभूत करें कि
``मदिरा पीता है'' केवल उन वस्तुओं के लिए सत्य है जो विधेय
$\Atom{\Obj{French}}{x}$ को संतुष्ट करती हैं, इत्यादि। इन रूढ़ियों और
अभिगृहीतों के साथ यह दिखाना कठिन नहीं है कि बहु-वर्गीय !!{sentence}s का
अनुवाद प्रथम-कोटि !!{sentence}s में होता है और बहु-वर्गीय !!{derivation}s
अनूदित होकर प्रथम-कोटि !!{derivation}s बनती हैं। बहु-वर्गीय !!{structure}s
और संगत प्रथम-कोटि !!{structure}s भी परस्पर ``अनूदित'' होती हैं; अतः
बहु-वर्गीय तर्कशास्त्र के लिए भी हमारे पास पूर्णता प्रमेय है।

\end{document}

